\section{The old integer margin absorbs the new index bound}
\label{sec:absorbed-bound}

Write \(C_*=L+\delta(A-1)\) for the new Pell coordinate and
retain the original source notation
\[
 Y=sM,\qquad p=2wr^2s^2M^3,\qquad
 \eta=k^2-4(c-kY)^2.
\]
Delete the separate inequality \(C_*^2<c^2\), and replace the
original requirement \(\eta>0\) by
\begin{equation}\label{eq:absorbed-margin}
                  \eta-C_*^2>0.
\end{equation}
This replacement needs no new coordinate: the square of the existing
coordinate \(C_*\) is already a quadratic monomial.

\begin{proposition}[Soundness of the absorbed bound]
Every rational solution of the modified source still satisfies the
old small-Pell strict inequality and the deleted index bound.
\end{proposition}
\begin{proof}
Equation~\eqref{eq:absorbed-margin} implies \(\eta>0\). Hence every
argument in the retained signed source that uses this inequality
remains valid. In particular, the half-trace square still excludes
\(s=0\) in the same noncircular order; this exclusion uses only
the original \(\eta>0\) consequence. Once \(s\ne0\), its
integrality and \(M\ge16\) give \(|Y|\ge16\). Since \(k\ne0\),
the old inequality then implies
\[
 \left|\frac{c}{k}-Y\right|<\frac12,
 \qquad |c|>|k|\left(|Y|-\frac12\right)>|k|.
\]
The strengthened margin also gives
\[
                 C_*^2<k^2-4(c-kY)^2\le k^2<c^2.
\]
Thus the deleted bound is available before the auxiliary Pell index
is identified. There is no dependence on knowing \(q=b^L\).
\end{proof}

\begin{proposition}[Completeness of the absorbed bound]
Every corrected positive source witness in the retained cofinite
instance range has a witness satisfying
\eqref{eq:absorbed-margin}, with the canonical auxiliary power
coordinate \(C_*=Y_L(A)\). The positivity coordinates can be chosen
anew; all source and Pell coordinates are retained.
\end{proposition}
\begin{proof}
The established positive reconstruction gives
\[
 w,s,k>0,\qquad k=Y_{r+1}(p),\qquad
 r\ge M\ge b^{128}>L,
 \qquad 0<A<5b^5s.
\]
The weaker lower bound \(M\ge b^{128}\) is sufficient here; the
canonical equality \(q=b^L\) provides a stronger bound. Since
\(w,r\ge1\), we obtain
\[
 p=2wr^2s^2M^3\ge2s^2M^3>100b^{10}s^2>4A^2.
\]
The strict middle inequality follows from
\(M^3\ge b^{384}>50b^{10}\), already true for \(b\ge2\).
For a positive Pell parameter the elementary recurrence bounds give
\[
 \begin{split}
 C_*^2=Y_L(A)^2
   &<(2A)^{2L-2}<p^{L-1}\\
   &<p^r<Y_{r+1}(p)=k.
 \end{split}
\]
Here \(L=5^{60}>2\), \(p\ge2\), and \(r>L-1\), so every
displayed strict inequality has its stated range.

It remains to compare this bound with the available margin. Put
\(e_0=c-kY\in\Z\). The original inequality is strict, so
\(k-2|e_0|\) is a positive integer. Consequently
\[
 \eta=(k-2|e_0|)(k+2|e_0|)\ge k>C_*^2.
\]
Thus \eqref{eq:absorbed-margin} holds. Its value is a positive
integer, so the same three triangular-number coordinates used by the
original positivity interface represent it. No additional integer
predicate or positivity variable is required.
\end{proof}

\paragraph{Exact row change.}
The retained positivity residual has the form
\[
       2\eta-2-\sum_{j=1}^3 u_j(u_j+1)=0.
\]
Subtracting the single quadratic monomial \(2C_*^2\) gives the
new residual. Delete the three private coordinates of the old
\(C_*^2<c^2\) wrapper and its single residual. This saves precisely
three coordinates and one quadratic residual, independently of the
three-coordinate saving from the signed-\(q\) argument.
